
% !TeX root = ../main.tex

%*****************************************
\chapter{Implementation}
\label{ch:implementation}
%*****************************************

\hint{This chapter should describe the details of the implementation addressing the following questions: \\ \\
1. What are the design decisions made? \\
2. What is the environment the approach is developed in? \\
3. How are components mapped to classes of the source code? \\
4. How do the components interact with each other?  \\
5. What are limitations of the implementation? }
\section{Design Decisions}
\label{sec:impl-design}

This section records the implementation-level decisions that shape the source code under
\texttt{code/}. The conceptual rationale for the overall approach is developed in
Chapter~\ref{ch:design}; here the focus is on the concrete choices made when translating
that design into executable PyTorch, and on the trade-offs that motivated each choice.
Where a decision follows an established method from the literature this is stated
explicitly; otherwise the decision is specific to this project.

\subsection{Three-Stage Pipeline Structure}
\label{sec:impl-design-pipeline}

The implementation is organised as three sequential stages, each with its own runner,
engine, and model class, so that a stage can be re-executed without re-running the stages
that precede it. Stage~1 is not a training stage in this codebase: it supplies frozen
pretrained encoder weights from the VideoMAE ViT-Base family~\cite{tong2022videomae},
loaded through \texttt{core/multimae.py} and selected by the \texttt{pretrained\_encoder}
configuration key. The two candidate corpora (K400 and SSV2) are compared only as initial
weights; the encoder is not retrained in Stage~1. Stage~2 trains the
\texttt{MultiModalMAE} model (\texttt{core/multimae.py}) via the runner
\texttt{runners/run\_pretrain.py} and the training loop in \texttt{engines/pretrain.py}.
Stage~3 trains the \texttt{MultiModalWaveformRegressor} (\texttt{core/waveform\_model.py})
via \texttt{runners/run\_waveform.py}; the Stage-3 training loop is implemented directly in
the runner rather than in a separate engine module. A diagnostic AU-occurrence probe
(\texttt{core/au\_probe.py}) is provided as an add-on that reads Stage-2 encoder weights
without modifying them.

\subsection{Asymmetric Multi-Modal Masking}
\label{sec:impl-design-amm}

The central masking decision is that visual and physiological streams are masked
asymmetrically. Visual streams (RGB and thermal-infrared) are masked so that at most 25\%
of their tokens are kept visible, leaving a large masked region for the decoder to
reconstruct, following the masked-reconstruction objective of
MAE~\cite{he2021maskedMAE} and VideoMAE~\cite{tong2022videomae}. Physiological streams are
masked without such a cap, so that a stream may be fully masked (a kept ratio of zero) when
the random draw requires it. The per-stream kept ratios are drawn from a single Dirichlet
distribution per optimisation step and shared across the whole batch, so that all samples
in a batch are masked with the same per-stream budget. A co-existence constraint requires
that at least one visual and at least one physiological stream be active in every batch;
this is enforced at construction time, so that a stream list containing only one kind of
stream is rejected rather than silently repaired. A further decision is that the
ViT~encoder~\cite{dosovitskiy2021image} operates on visible tokens only: no mask
placeholder tokens are inserted into the encoder input, which removes the quadratic
attention cost~\cite{vaswani2017attention} of attending over masked positions. Mask tokens
are re-inserted only for the decoder.

\subsection{Span Masking for Physiological Streams}
\label{sec:impl-design-span}

For one-dimensional physiological signals, the default masking is random per-token masking,
which preserves bit-for-bit agreement with earlier runs. An opt-in span-masking mode is
also provided. The motivation is that scattered random masking at moderate ratios can be
solved by temporal interpolation of the visible signal alone, without using the video, which
would let the model learn a shortcut that does not transfer to the sensor-failure setting.
Masking contiguous spans removes the surrounding visible context that interpolation relies
on. The span length and the number of spans are drawn once per step and shared across the
batch; this is a load-bearing invariant, because the encoder consumes the batch-maximum of
the per-sample mask, and a per-sample span draw would leak masked tokens into the encoder.

\subsection{Model Registry and Single Entry Point}
\label{sec:impl-design-registry}

All models are constructed through a single factory function, \texttt{create\_model} in
\texttt{models/build.py}, which dispatches on a registered model name. Model classes
register themselves through a \texttt{@register\_model} decorator in
\texttt{core/registry.py}. Stage~3 and the AU probe use dedicated builder functions rather
than the generic factory. This keeps model construction in one place and makes the set of
available models explicit.

\subsection{Configuration System}
\label{sec:impl-design-config}

Experiments are configured with YAML files that supply \emph{defaults only} to the
\texttt{argparse} parser. The precedence order is: an explicit command-line flag overrides
a YAML value, which in turn overrides the hard-coded default. This means a YAML file is a
named starting point rather than a fixed recipe, and any single run can be varied without
editing the file. The same mechanism is shared by all runners through a common helper in
\texttt{runners/\_common.py}.

\subsection{Reproducibility: Distributed Setup, Seeding, and Precision}
\label{sec:impl-design-repro}

Distributed bootstrap is centralised in \texttt{utils/dist.py}; each runner calls a common
\texttt{init\_env} helper that initialises the NCCL backend and the device. Random seeding
is likewise centralised: the base seed is offset by the per-process rank, so that each
distributed rank draws a distinct but deterministic stream. Mixed precision is opt-in,
selected by an \texttt{amp\_dtype} flag whose default (\texttt{none}) runs in full
32-bit floating point and is bit-identical to the unflagged path; when enabled, a custom
loss scaler in \texttt{utils/native\_scaler.py} performs loss scaling and returns the
gradient norm.

\subsection{Normalisation Ownership}
\label{sec:impl-design-norm}

A deliberate decision is to place label normalisation in the dataset rather than in the
model. For the region-of-interest data lineages, the dataset applies a per-clip z-score to
the physiological target, and the model is pinned to an identity normalisation so that it
never re-normalises. The retired normalisation options are explicitly refused by a check
function rather than silently ignored, so that a stale configuration fails loudly.

\subsection{Subject-Disjoint Cross-Validation}
\label{sec:impl-design-cv}

Evaluation uses a subject-disjoint K-fold scheme: subjects are partitioned once into
folds, and each fold in turn serves as the test set while the remaining folds form the
training set, so that no subject appears in both the training and test sets of any fold.
A small validation split is optionally carved from the training folds with a fixed
per-fold seed. The partition rule (seed and optional stratification by sex) is recorded in
a manifest file so that Stage~2 and Stage~3 can be run with the identical split.

\subsection{Attention Implementation}
\label{sec:impl-design-attention}

Self-attention is implemented with the fused scaled-dot-product-attention kernel rather
than a hand-written attention. At the full 224-pixel input resolution the token count is
large enough that the naive path is impractical in memory, so the fused kernel is a
requirement at that resolution rather than an optimisation.

\subsection{Stage-2 / Stage-3 Geometry Lock}
\label{sec:impl-design-geometry}

Because Stage~3 transfers Stage-2 encoder weights verbatim, the two stages must operate on
identical token geometry. A check at load time verifies that the tubelet size, input
resolution, and frame count match, and that the signal window per token is consistent with
the video temporal stride; a mismatch raises an error at load rather than producing a
silently misaligned model.

\subsection{Metric Logging Discipline}
\label{sec:impl-design-logging}

Training metrics are recorded through a single metric-logger utility in
\texttt{utils/logger.py}. A standing rule is that tensor values are converted to Python
scalars before being stored, so that the autograd graph is not retained in GPU memory; this
prevents out-of-memory growth during long runs.

\section{Architecture}

\section{Interaction of Components}

\section{Summary}